\chapter{测度的积分}

本章中，T 表示一个局部紧空间，\(\mu\) 表示 T 上一个正测度。对集 E 的每个子集 A，\(\varphi_{A}\) 表示 A 的特征函数（若不致混淆）。所谓数值函数，我们总指取值于 \(\overline{R}\) 的函数，因而可能取值 \(+\infty\) 与 \(-\infty\) 。定义在 E 上的正数值函数所成之集记作 \(\mathcal{F}_{+}(E)\) ，若不致混淆也简记作 \(F_{+}\) 。我们约定把乘积 \(0 \cdot (+\infty)\) 与 \(0 \cdot (-\infty)\) 定义为取值 0；于是，若 f 是定义在 E 上的数值函数而 A 是 E 的子集，则 \(f\varphi_{A}\) 表示在 A 上与 f 相同、在 CA 上等于 0 的函数。对局部紧空间的每个点 a，\(\varepsilon_{a}\) 表示在点 a 放置单位质量所定义的测度（第三章，§1 第 3 号）。

本质上积分（相应地，本质可积函数）的概念将在 §1 中定义；我们将看到，当局部紧空间 T 在无穷远处可数时（GT, I, §9 第 9 号），它与上积分（相应地，可积函数）的概念一致。

因此，只关心在无穷远处可数的局部紧空间上的积分的读者可以略去 §1 的第 1 至 3 号不读；在本章其余部分，当所考虑的空间在无穷远处可数时，只要删去“本质的”与“本质上”等词，并把记号 \(\mu^{\bullet}\) 换成 \(\mu^{*}\) 、记号 \(\int^{\bullet}\) 换成 \(\int^{*}\) ，就得到有效的命题。

在 §1 至 §4 中，测度一词总是指正测度；其他测度将视情形明确地称为未必正的测度或复测度。

